\documentclass[11pt]{article}
\usepackage[T1]{fontenc}
\usepackage{lmodern,amsmath,amssymb,amsthm,geometry,microtype}
\usepackage[hyphens]{url}
\usepackage[colorlinks=true,linkcolor=blue,citecolor=blue,urlcolor=blue]{hyperref}
\geometry{margin=1in}
\newtheorem{theorem}{Theorem}
\newtheorem{corollary}{Corollary}
\newcommand{\E}{\mathbb E}
\newcommand{\Var}{\operatorname{Var}}
\newcommand{\dist}{\operatorname{dist}}
\hypersetup{pdftitle={A spectral remainder in the Hellinger-to-information implication},pdfauthor={Tom Ayre},pdfsubject={A quantitative refinement conditional on Hellinger contraction}}
\title{A spectral remainder in the\\Hellinger-to-information implication}
\author{Tom Ayre\\\small\href{mailto:ayre.tom@gmail.com}{ayre.tom@gmail.com}}
\date{October 10, 2026}
\begin{document}
\maketitle

\begin{abstract}
We retain a quadratic remainder in the standard implication from the
Boolean Hellinger inequality to the Courtade--Kumar information bound.
Assuming the Hellinger inequality for a Boolean function $f$, its
information deficit is at least
$[\rho^2-\Var(T_\rho f)]/6$. Combined with the
Friedgut--Kalai--Naor theorem, this gives distance to a signed coordinate
at most $C\Delta_\rho(f)/[\rho^2(1-\rho^2)]$ for a universal constant $C$.
The general-function statement depends on the full-bias Hellinger
inequality asserted in the cited source. For antipodal functions, the
same conclusion follows from Ky's separate Hellinger theorem.
This is a potentially new quantitative refinement of an existing
argument; publication priority has not been established.
\end{abstract}

\section{Setting and source assumption}
Let $n\ge1$, let $X$ be uniform on $\{-1,1\}^n$, and let
$f:\{-1,1\}^n\to\{-1,1\}$. For $0\le\rho\le1$, obtain $Y_\rho$
by independently flipping each coordinate of $X$ with probability
$(1-\rho)/2$. The noise operator is
\[
 T_\rho f(y)=\E[f(X)\mid Y_\rho=y].
\]
All expectations use the uniform cube unless otherwise stated.
All logarithms are natural, and mutual information is measured in nats.
With $0\log0=0$, write
\[
 h_2(p)=-p\log p-(1-p)\log(1-p),\quad
 \psi(r)=\log2-h_2((1+r)/2),\quad R(r)=\sqrt{1-r^2}.
\]
The information deficit from a signed coordinate is
\[
 \Delta_\rho(f)=\psi(\rho)-I(f(X);Y_\rho).
\]
For $S\subseteq[n]$, put
$\widehat f(S)=\E[f(X)\prod_{i\in S}X_i]$ and
$W_1(f)=\sum_{i=1}^n\widehat f(\{i\})^2$.

The input to the argument is the Hellinger inequality
\begin{equation}\label{eq:H}
 R(\E f)-\E R(T_\rho f)\le1-R(\rho).
 \tag{H}
\end{equation}
Theorem 1.1 of the repository manuscript \cite{source} asserts
\eqref{eq:H} for all such $f$, including functions with nonzero mean.
We prove an implication from this input and do not give a new proof
of the general Hellinger inequality. The qualitative implication to
mutual information is due to Anantharam et al.\ \cite{ABCJN}.

\section{The spectral remainder}
\begin{theorem}\label{thm:main}
Let $f:\{-1,1\}^n\to\{-1,1\}$ and $0\le\rho\le1$.
If \eqref{eq:H} holds for this $f$ and $\rho$, then
\begin{align}
 \Delta_\rho(f)
 &\ge\frac16\left(\rho^2-\Var(T_\rho f)\right)\label{eq:main}\\
 &=\frac{\rho^2}{6}\left((\E f)^2+
 \sum_{|S|\ge2}(1-\rho^{2|S|-2})\widehat f(S)^2\right).\notag
\end{align}
Consequently,
\begin{equation}\label{eq:level-one}
 \Delta_\rho(f)\ge\frac{\rho^2(1-\rho^2)}6\bigl(1-W_1(f)\bigr).
\end{equation}
\end{theorem}
\begin{proof}
Define $\phi(t)=h_2((1+\sqrt{1-t^2})/2)$ for $0\le t\le1$.
For $0<t<1$, putting $v=\sqrt{1-t^2}$ gives
\[
 \phi'(t)=\frac{t\operatorname{atanh}(v)}v,\qquad
 \phi''(t)=\frac{\operatorname{atanh}(v)-v}{v^3}
 =\sum_{j\ge1}\frac{v^{2j-2}}{2j+1}\ge\frac13.
\]
Moreover $\phi'(t)\to0$ as $t\downarrow0$. It follows that
$q(t)=\phi(t)-t^2/6$ is convex and nondecreasing on $[0,1]$.

Put $g=T_\rho f$, $m=\E f$, $Z=R(g)$,
$h=R(m)$, $b=\E Z$, and $s=R(\rho)$.
Concavity of $R$ and $\E g=m$ give $0\le b\le h\le1$.
Assumption \eqref{eq:H} gives $h-b\le1-s$.
Jensen's inequality and convexity therefore yield
\begin{equation}\label{eq:transfer}
 q(h)-\E q(Z)\le q(h)-q(b)
 \le q(1)-q(1-h+b)\le q(1)-q(s).
\end{equation}
For the middle inequality, translate the interval $[b,h]$ rightward
by $1-h$; increments of a convex function over an interval of fixed
length increase with its left endpoint. The last inequality follows
from $1-h+b\ge s$ and monotonicity of $q$.

The posterior-mean identity gives
$I(f(X);Y_\rho)=\phi(h)-\E\phi(Z)$, and
$\psi(\rho)=\phi(1)-\phi(s)$. Substituting
$q(t)=\phi(t)-t^2/6$ into \eqref{eq:transfer} gives
\[
 \Delta_\rho(f)\ge\frac16(1-s^2-h^2+\E Z^2)
 =\frac16(\rho^2+m^2-\E g^2),
\]
which is \eqref{eq:main}. Parseval and
$\widehat{T_\rho f}(S)=\rho^{|S|}\widehat f(S)$ give its Fourier
expansion. For $|S|\ge2$, $1-\rho^{2|S|-2}\ge1-\rho^2$;
also $1\ge1-\rho^2$ for the mean term. This proves
\eqref{eq:level-one}. The argument includes the endpoints by continuity.
\end{proof}

\begin{corollary}[Stability]\label{cor:stability}
There is a universal $C<\infty$ such that, if \eqref{eq:H} holds for
$f$ and $0<\rho<1$, then
\[
 \dist(f,\mathrm{dictators}):=
 \min_{\substack{1\le i\le n\\\sigma\in\{-1,1\}}}
 \Pr[f(X)\ne\sigma X_i]
 \le\frac{C\Delta_\rho(f)}{\rho^2(1-\rho^2)}.
\]
\end{corollary}
\begin{proof}
Set $\eta=1-W_1(f)$. The Friedgut--Kalai--Naor theorem \cite{FKN}
gives a Boolean function $a$ depending on at most one coordinate,
with $\Pr[f\ne a]\le K\sum_{|S|\ge2}\widehat f(S)^2\le K\eta$,
for some universal $K\ge1$.
If $\eta<\min(1/4,1/(8K))$, then $(\E f)^2\le\eta$ implies
$|\E f|<1/2$. Every constant consequently has disagreement greater
than $1/4$ with $f$, whereas the FKN approximant has disagreement less
than $1/8$. Thus $a$ is a signed coordinate.
For larger $\eta$, the trivial bound
$\dist(f,\mathrm{dictators})\le1/2$ gives the same conclusion
$\dist(f,\mathrm{dictators})\le C_0\eta$ after enlarging the universal
constant. Apply \eqref{eq:level-one} with $C=6C_0$.
\end{proof}

\begin{corollary}[Antipodal functions]\label{cor:antipodal}
If $f(-x)=-f(x)$ for every $x$, then the conclusions of
Theorem~\ref{thm:main} and Corollary~\ref{cor:stability} follow from
Ky's antipodal Hellinger theorem \cite[Theorem 1.1]{Ky}, without using
the general full-bias input of \cite{source}.
\end{corollary}
\begin{proof}
Antipodality gives $\E f=0$. Ky's theorem states
$\E\sqrt{1-(T_\rho f)^2}\ge\sqrt{1-\rho^2}$, which is exactly
\eqref{eq:H} when $\E f=0$. Apply the preceding results.
\end{proof}

\section{Comparison and scope}
Ky and Tran \cite[Theorem 1.1]{KT} prove
\[
 \dist(f,\mathrm{dictators})\le
 \frac{10^8\Delta_\rho(f)}{\rho^2(1-\rho)h_2((1-\rho)/2)}.
\]
For increasing functions they remove $1-\rho$ from that denominator.
Their Section 5.6 leaves open whether this extra factor can be removed
for general functions. Conditional on \eqref{eq:H}, our denominator
has order $1-\rho$ as $\rho\uparrow1$, improving the general bound's
order $(1-\rho)^2\log(1/(1-\rho))$. Our unspecified universal constant
means this is an asymptotic comparison, not a pointwise comparison of
numerical constants. We do not achieve the increasing-function
order $(1-\rho)\log(1/(1-\rho))$ or resolve the sharp general question.

Dimension-independent stability is already known. The potential
contribution here is the explicit spectral remainder and the resulting
noise dependence. The convex-profile transfer is already in
\cite{ABCJN} and is reproduced in \cite{source}. Our extra step is to
subtract $t^2/6$ before that transfer. This coefficient is the largest
allowed by this global quadratic-subtraction argument, since
$\phi''(t)\to1/3$ as $t\uparrow1$. We do not claim that $1/6$ is
the optimal coefficient in \eqref{eq:main}.

The separate antipodal input \cite{Ky} was posted on October 7, 2026.
Its Section 6.2 discusses quantitative stability as a further question.
Corollary~\ref{cor:antipodal} supplies an information-stability
consequence; it does not quantify all losses in that paper's matrix
argument or give its equality classification for general product channels.
A targeted literature check has not established priority for the present
refinement. This manuscript is a research note, not a claim of a new
construction or a new proof of the general Hellinger inequality.

\section{Checks and disclosure}
The accompanying script \path{check_hellinger_stability.py} enumerates
every Boolean function in dimensions one through four at eight
correlations. All 526,496 cases passed checks of the Hellinger input,
the spectral remainder, and the Fourier identity, at floating-point
tolerance $10^{-12}$. These finite diagnostics are not a proof.

The source's complete \path{verification/verify_all.py} workflow was
also rerun on October 7, 2026 with Python 3.11.14, SymPy 1.14.0 and
\texttt{clang++}; all supplied arithmetic-certificate stages passed.
The accompanying \path{source-verification-summary.json} records the
source revision, file hashes, and outcomes. This checks finite
certificates, not the entire analytic proof of \cite{source}.
The note, diagnostic code, and internal reviews were prepared with
assistance from Codex. No Lean kernel verification or external peer
review is claimed.

\begin{thebibliography}{9}
\bibitem{source} OpenAI, \emph{Hellinger contraction with arbitrary Boolean
output bias}, September 24, 2026, Theorem 1.1. In \texttt{openai/math},
revision \texttt{adc7f1241b42e322a6451854ab7e4b4c146bf78a}.
\href{https://github.com/openai/math/blob/adc7f1241b42e322a6451854ab7e4b4c146bf78a/preprints/Hellinger-contraction-with-arbitrary-Boolean-output-bias-September-24-2026/main.pdf}{Manuscript and versioned repository source}.
\bibitem{ABCJN} V. Anantharam, A. Bogdanov, A. Chakrabarti, T. S. Jayram,
and C. Nair, \emph{A conjecture regarding optimality of the dictator
function under Hellinger distance}, ITA workshop manuscript, 2017,
Proposition 1 and Lemma 1.
\url{https://chandra.ie.cuhk.edu.hk/pub/papers/HC/hel-conj.pdf}.
\bibitem{FKN} E. Friedgut, G. Kalai, and A. Naor, \emph{Boolean functions
whose Fourier transform is concentrated on the first two levels},
Advances in Applied Mathematics 29 (2002), 427--437.
\bibitem{KT} V. K. Ky and T. Tran, \emph{Stability of the Courtade--Kumar
inequality}, arXiv:2609.31797v1, September 25, 2026.
\url{https://arxiv.org/abs/2609.31797v1}.
\bibitem{Ky} V. K. Ky, \emph{The Hellinger conjecture for antipodal Boolean
functions}, arXiv:2610.10618v1, October 7, 2026.
\url{https://arxiv.org/abs/2610.10618v1}.
\end{thebibliography}
\end{document}
